% ============================================================
%  Active Subgraph Skip Profiles — Research Group Slides
%  Data source: graphs_0505  (generated 2025-05-05)
%  Models: GPT-2 (12L$\times$12H)  |  GPT-2-medium (24L$\times$16H)
%  Compile:  pdflatex slides_0505.tex
% ============================================================
\documentclass[aspectratio=169,10pt]{beamer}

\usetheme{Madrid}
\usecolortheme{default}

% ── Custom palette (matches the visualisation colour scheme) ──────────────
\definecolor{teal}   {RGB}{23,165,137}   % residual/skip arcs
\definecolor{attnred}{RGB}{214, 39, 40}  % active attention head
\definecolor{mlpblue}{RGB}{ 31,119,180}  % active FFN block
\definecolor{darkgray}{RGB}{80,80,80}
\definecolor{lightbg} {RGB}{245,245,245}

\setbeamercolor{frametitle}{bg=teal,fg=white}
\setbeamercolor{title}{bg=teal!90!black,fg=white}
\setbeamercolor{structure}{fg=teal}
\setbeamercolor{block title}{bg=teal!20,fg=teal!80!black}
\setbeamercolor{block body}{bg=lightbg}
\setbeamercolor{alerted text}{fg=attnred}

% ── Packages ──────────────────────────────────────────────────────────────
\usepackage[utf8]{inputenc}
\usepackage{pifont}
\newcommand{\cmark}{\ding{51}}
\newcommand{\xmark}{\ding{55}}
\usepackage{booktabs}
\usepackage{multirow}
\usepackage{colortbl}
\usepackage{xcolor}
\usepackage{array}
\usepackage{graphicx}
\graphicspath{{../graphs_0505/}}   % PNG files live one level up
\usepackage{tikz}
\usetikzlibrary{positioning,arrows.meta,calc}

% convenience: active / inactive cell macros
\newcommand{\yes}[1]{\textcolor{teal}{\textbf{#1}}}
\newcommand{\no}[1] {\textcolor{gray}{#1}}
\newcommand{\hi}[1] {\textcolor{attnred}{\textbf{#1}}}

% ── Title metadata ────────────────────────────────────────────────────────
\title[Skip Profiles in GPT-2 Family]{Active Subgraph Skip Profiles:\\
       Layer-level Computation Routing in GPT-2 Family Models}
\subtitle{graphs\_0505 analysis — reasoning suite}
\author{LLM Path Analysis Project}
\date{May 2025}

% ─────────────────────────────────────────────────────────────────────────
\begin{document}
% ─────────────────────────────────────────────────────────────────────────

% ══════════════════════════════════════════════════════════════════════════
\begin{frame}
\titlepage
\end{frame}

% ══════════════════════════════════════════════════════════════════════════
\begin{frame}{Slide 1 — Experiment Setup}

\begin{columns}[T]

%% LEFT ─ models & tasks
\begin{column}{0.46\textwidth}
\begin{block}{Models}
\footnotesize
\begin{tabular}{llcc}
\toprule
Model & Layers & Heads & Params \\
\midrule
\texttt{gpt2}        & 12 & 12 & 117M \\
\texttt{gpt2-medium} & 24 & 16 & 345M \\
\bottomrule
\end{tabular}\\[4pt]
Both run on CPU, \texttt{eval()} mode.
\end{block}

\vspace{4pt}
\begin{block}{Task suite: \texttt{reasoning} (6 tasks)}
\footnotesize
\begin{tabular}{ll}
\toprule
Label & Prompt (truncated) \\
\midrule
1-hop chain  & \emph{Alice is the mother of Bob. Alice's child is} \\
2-hop chain  & \emph{Alice $\to$ Bob $\to$ Carol. Alice's grandchild is} \\
3-hop chain  & \emph{Alice $\to$ Bob $\to$ Carol $\to$ Dana. \dots grandchild is} \\
Cat.\ syllogism & \emph{All birds have wings. Penguins are birds. So\dots} \\
Negation+ded.  & \emph{No reptiles are warm-blooded. All mammals\dots} \\
Counterfactual & \emph{In a world where cats bark\dots barking is a} \\
\bottomrule
\end{tabular}
\end{block}
\end{column}

%% RIGHT ─ methodology
\begin{column}{0.50\textwidth}
\begin{block}{Active/inactive criterion (new)}
\footnotesize
Each block (attention or FFN) is classified by its\\
\textbf{output-to-input magnitude ratio}:
\[
r_{\text{attn}}^{(l)} = \frac{\|\mathtt{attn\_out}^{(l)}\|_2}{\|\mathtt{resid\_pre}^{(l)}\|_2},
\quad
r_{\text{ffn}}^{(l)} = \frac{\|\mathtt{mlp\_out}^{(l)}\|_2}{\|\mathtt{resid\_mid}^{(l)}\|_2}
\]
\textbf{Inactive} (skip-dominant): $r < 0.05$ \quad
\textcolor{teal}{$\boldsymbol{\Rightarrow}$ thick teal U-arc}\\
\textbf{Active} (block transforms): $r \geq 0.05$ \quad
\textcolor{gray}{$\Rightarrow$ thin dashed arc}
\end{block}

\vspace{4pt}
\begin{block}{Attribution scoring (AtP)}
\footnotesize
Per-head score at \texttt{hook\_z}:
\[
s(l,h) = \textstyle\text{mean}_{t,d}\,|\,\nabla_z \log p \odot z\,|
\]
Nucleus coverage $\rho=0.90$: keep fewest edges\\
covering 90\% of total attribution mass.\\
Model params frozen (\texttt{requires\_grad\_(False)})\\
before backward to avoid $O(P)$ gradient cost.
\end{block}
\end{column}

\end{columns}
\end{frame}

% ══════════════════════════════════════════════════════════════════════════
\begin{frame}{Slide 2 — Finding 1: Opposing First-Layer Strategies}

\begin{columns}[T]

\begin{column}{0.48\textwidth}
\begin{block}{\textcolor{attnred}{GPT-2} — attention-dominant entry}
\footnotesize
\begin{tabular}{ccccc}
\toprule
Layer & \multicolumn{2}{c}{Attn ratio $r$} & \multicolumn{2}{c}{FFN ratio $r$} \\
      & status & value & status & value \\
\midrule
L0 & \yes{\cmark{}} & \hi{5.2–5.9} & \yes{\cmark{}} & \hi{1.2–1.8} \\
L1 & \yes{\cmark{}} & 0.17         & \yes{\cmark{}} & 0.43–0.71 \\
L2 & \yes{\cmark{}} & 0.08–0.11    & \yes{\cmark{}} & 1.12–1.84 \\
L3 & varies  & 0.03–0.07    & \yes{\cmark{}} & 0.11–0.13 \\
\bottomrule
\end{tabular}

\medskip
\textbf{L0 attention $r \approx 5.5$}: attn output is\\
$5\times$ the stream magnitude — heads dominate\\
the very first transformation.\\[4pt]
L2 FFN is a secondary spike ($r \approx 1.5$), the\\
first ``feature extraction'' checkpoint.
\end{block}
\end{column}

\begin{column}{0.48\textwidth}
\begin{block}{\textcolor{mlpblue}{GPT-2-medium} — FFN-dominant entry}
\footnotesize
\begin{tabular}{ccccc}
\toprule
Layer & \multicolumn{2}{c}{Attn ratio $r$} & \multicolumn{2}{c}{FFN ratio $r$} \\
      & status & value & status & value \\
\midrule
L0 & \yes{\cmark{}} & 0.34–0.36    & \yes{\cmark{}} & \hi{36–38} \\
L1 & \yes{\cmark{}} & 0.14–0.16    & \yes{\cmark{}} & 0.19–0.21 \\
L2 & \yes{\cmark{}} & 0.11–0.15    & \yes{\cmark{}} & 0.19–0.28 \\
L3 & \yes{\cmark{}} & 0.06–0.07    & \yes{\cmark{}} & 0.86–1.71 \\
\bottomrule
\end{tabular}

\medskip
\textbf{L0 FFN $r \approx 37$}: FFN output is\\
$37\times$ the stream magnitude — the first feed-\\
forward layer essentially rewrites the token.\\[4pt]
Attention at L0 is present ($r\approx0.35$) but\\
plays a minor role compared to the FFN.
\end{block}
\end{column}

\end{columns}

\vspace{6pt}
\begin{alertblock}{Interpretation}
\footnotesize
Both models always activate L0, but which \emph{sub-block} dominates differs entirely.
GPT-2 is \textbf{attention-first}; GPT-2-medium is \textbf{FFN-first}.
This suggests the two training runs converge on different strategies for converting token embeddings into
distributional representations — a difference invisible to loss curves but visible in magnitude profiles.
\end{alertblock}

\end{frame}

% ══════════════════════════════════════════════════════════════════════════
\begin{frame}{Slide 3 — Finding 2: Attention Silence Zones and the Bilateral Dead Zone}

\begin{columns}[T]

\begin{column}{0.52\textwidth}
\begin{block}{Attention-silent layers per task}
\footnotesize
\setlength{\tabcolsep}{4pt}
\begin{tabular}{lcc}
\toprule
Task & GPT-2 attn $\times$ layers & GPT-2-med attn $\times$ layers \\
\midrule
1-hop chain  & \no{L3–L9 (7 layers)} & \no{L4–L20 (17 layers)} \\
2-hop chain  & \no{L3–L4 (2 layers)} & \no{L4–L16 (13 layers)} \\
3-hop chain  & \yes{none}            & \yes{none} \\
Syllogism    & \no{L3–L7 (5 layers)} & \no{L4–L14 (11 layers)} \\
Negation     & \no{L3–L4, L6}        & \no{L4–L16 (13 layers)} \\
Counterfactual & \no{L4 only}        & \no{L4–L14, L16, L20} \\
\bottomrule
\end{tabular}

\medskip
\textbf{GPT-2}: FFN is \emph{never} inactive ($r \geq 0.07$ always).\\
Middle silence = attention off, FFN-only highway.\\[4pt]
\textbf{GPT-2-medium, 1-hop}: L7–L13 show \textbf{both}\\
attn \emph{and} FFN inactive simultaneously ($r < 0.05$).\\
The residual stream coasts 7 layers \textbf{untouched}.
\end{block}
\end{column}

\begin{column}{0.44\textwidth}
\begin{block}{3-hop chain: no silence}
\footnotesize
The 3-hop chain is the \emph{only} task in GPT-2-medium\\
where every layer shows attn \yes{\cmark{}}.\\[6pt]
This is \textbf{counter-intuitive}: the hardest chain task\\
uses the \emph{most} layers, while the easiest (1-hop)\\
uses the fewest — concentrating computation\\
into $\approx$6 active layers then coasting.
\end{block}

\vspace{4pt}
\begin{block}{Silence length scales with task}
\footnotesize
As chain depth increases, the attention-silent zone\\
\textbf{shrinks}: the model recruits progressively more\\
middle-layer attention heads to resolve each\\
additional relational hop.\\[6pt]
$\rightarrow$ Silence length is a proxy for\\
\quad task complexity in chain reasoning.
\end{block}

\vspace{4pt}
\begin{alertblock}{\footnotesize Key contrast with prior view}
\footnotesize
The "later layers refine" hypothesis predicts complex\\
tasks use \emph{more} late-layer processing.\\
Here: complex chains wake \emph{earlier} and keep\\
\emph{more} layers active throughout — not just at the end.
\end{alertblock}
\end{column}

\end{columns}
\end{frame}

% ══════════════════════════════════════════════════════════════════════════
\begin{frame}{Slide 4 — Finding 3 \& 4: Early Spike Decay + Universal Output Integrator}

\begin{columns}[T]

%% LEFT
\begin{column}{0.50\textwidth}
\begin{block}{Finding 3 — Early FFN spike \emph{decreases} with task complexity}
\footnotesize
The early prominent FFN layer decays monotonically\\
as chain depth increases:

\medskip
\begin{tabular}{lcc}
\toprule
Task & GPT-2 L2 FFN $r$ & GPT-2-med L3 FFN $r$ \\
\midrule
1-hop     & \hi{1.84} & \hi{1.71} \\
2-hop     & 1.41      & 1.15 \\
3-hop     & 1.12      & 0.86 \\
Syllogism & 1.43      & 1.15 \\
Negation  & 1.42      & 1.13 \\
Counterfactual & 1.34 & 1.05 \\
\bottomrule
\end{tabular}

\medskip
Simple tasks front-load more energy into the early\\
FFN (acting as a one-shot lookup).
Complex tasks distribute this budget across layers.\\[4pt]
Logic tasks (syllogism, negation) cluster near the\\
2-hop level — consistent with 2-premise retrieval.
\end{block}
\end{column}

%% RIGHT
\begin{column}{0.46\textwidth}
\begin{block}{Finding 4 — Universal output integrator}
\footnotesize
The \emph{final} attention layer fires strongly for\\
\textbf{every task} in both models:

\medskip
\begin{tabular}{lcc}
\toprule
Task & GPT-2 L11 & GPT-2-med L23 \\
     & attn $r$  & attn $r$ \\
\midrule
1-hop     & 0.89 & 0.39 \\
2-hop     & 0.96 & 0.41 \\
3-hop     & \hi{0.99} & \hi{0.42} \\
Syllogism & 0.90 & 0.38 \\
Negation  & 0.86 & 0.41 \\
Counterfactual & 0.93 & 0.38 \\
\bottomrule
\end{tabular}

\medskip
GPT-2 L11: always \yes{\cmark{}}, ratio \textbf{increases} with\\
hop depth ($0.89 \to 0.99$).\\[4pt]
GPT-2-medium L23: always \yes{\cmark{}}, lower absolute\\
ratio because stream magnitude has grown over\\
24 layers. GPT-2-medium L22 FFN also\\
shows a consistent ``pre-final'' spike ($r \approx 0.30$).\\[4pt]
$\Rightarrow$ Both models learn a \textbf{dedicated last-layer}\\
\textbf{attention integrator}, task-invariant but\\
calibrated to chain depth.
\end{block}
\end{column}

\end{columns}
\end{frame}

% ══════════════════════════════════════════════════════════════════════════
\begin{frame}{Slide 5 — Cross-task Clustering and Summary}

\begin{columns}[T]

\begin{column}{0.46\textwidth}
\begin{block}{Task clustering by silence signature (GPT-2-medium)}
\footnotesize
\begin{tabular}{ll}
\toprule
Cluster & Tasks \\
\midrule
\textbf{Deep silence} (wake $\geq$ L21) & 1-hop chain \\
\textbf{Long silence} (wake L17)         & 2-hop chain, \hi{negation+ded.} \\
\textbf{Medium silence} (wake L15)       & Syllogism, \hi{counterfactual} \\
\textbf{No silence}                      & 3-hop chain (unique) \\
\bottomrule
\end{tabular}

\medskip
\textbf{Negation+deduction groups with 2-hop},\\
not with the other logic tasks.\\
Supports the hypothesis that the model routes\\
``No X / All Y / therefore\dots'' through the\\
same 2-premise lookup circuit as a 2-hop chain.\\[6pt]
\textbf{Counterfactual} is the least chain-like logic\\
task — its silence is shortest among logic tasks,\\
suggesting additional world-model access.
\end{block}
\end{column}

\begin{column}{0.50\textwidth}
\begin{block}{Summary of findings}
\footnotesize
\renewcommand{\arraystretch}{1.25}
\begin{tabular}{p{3.2cm}p{4.0cm}}
\toprule
Finding & Implication \\
\midrule
L0: attn-dominant (GPT-2) vs.\ FFN-dominant (GPT-2-med) &
Different tokenisation strategies; not visible from loss \\
\midrule
Attention-silent zones shrink with hop depth &
Silence length $\approx$ task complexity proxy \\
\midrule
GPT-2-med 1-hop: bilateral dead zone L7–L13 &
Residual stream coasts 7 layers untouched \\
\midrule
Early FFN spike $\downarrow$ as hop depth $\uparrow$ &
Complex tasks distribute, not concentrate, budget \\
\midrule
Final-layer attn $r$ scales with hop depth &
Universal output integrator, strength $\propto$ complexity \\
\midrule
Negation clusters with 2-hop & Two-premise circuit reuse \\
\bottomrule
\end{tabular}
\end{block}

\vspace{4pt}
\begin{block}{Open questions}
\footnotesize
\begin{itemize}
  \item Does the bilateral dead zone persist for $>$3-hop chains in GPT-2-medium?
  \item Is the L0 attn/FFN dominance flip consistent across the full Pythia family?
  \item Can interventional patching \emph{into} the silent zone change task routing?
  \item Do larger models (GPT-2-large, GPT-2-XL) continue the early-FFN-spike-decay trend?
\end{itemize}
\end{block}
\end{column}

\end{columns}

\vspace{4pt}
\begin{block}{\footnotesize Code \& data}
\footnotesize
\texttt{active\_subgraph\_dot.py} \quad
\texttt{--models gpt2 gpt2-medium --suite reasoning --mag\_threshold 0.05 --out graphs\_0505/compare}\\
Results: \texttt{graphs\_0505/compare\_\{gpt2,gpt2\_medium\}\_task\{0..5\}.\{md,dot,png,svg\}}
\end{block}

\end{frame}

% ══════════════════════════════════════════════════════════════════════════
%  Image slides — actual active subgraph outputs from graphs_0505
% ══════════════════════════════════════════════════════════════════════════

% ── Image slide 1: GPT-2 3-hop chain ─────────────────────────────────────
% Chosen because it has the best aspect ratio for slides (tall + narrow).
% Key feature: NO attention silence — every layer fires, unlike 1-hop.
\begin{frame}[plain]
\begin{tikzpicture}[remember picture, overlay]

  % Teal header bar
  \fill[teal] (current page.north west) rectangle
              ([yshift=-0.9cm]current page.north east);
  \node[anchor=west, text=white, font=\bfseries\footnotesize]
        at ([xshift=0.4cm, yshift=-0.45cm]current page.north west)
        {GPT-2 (12L$\times$12H) \quad 3-hop Chain \quad $k=18$ \quad
         \texttt{compare\_gpt2\_task2.png}};

  % ── Image (left column, full height) ─────────────────────────────────
  \node[anchor=north west] (img)
        at ([xshift=0.2cm,yshift=-1.0cm]current page.north west) {%
    \includegraphics[height=0.86\paperheight,
                     keepaspectratio]{compare_gpt2_task2.png}%
  };

  % ── Annotation panel (right of image) ────────────────────────────────
  \node[anchor=north west, text width=7.4cm,
        font=\footnotesize, align=left]
        at ([xshift=3.0cm,yshift=-1.1cm]current page.north west) {

    \textcolor{teal}{\textbf{What to look at}}

    \medskip
    \textbf{Skip arc colour} (left U-shape beside each layer):
    \begin{itemize}\setlength\itemsep{2pt}
      \item \textcolor{teal}{\textbf{Thick teal}} = block INACTIVE ($r<0.05$);\\
            skip is the dominant path
      \item \textcolor{gray}{Thin dashed} = block ACTIVE ($r\geq0.05$);\\
            block transforms the stream
    \end{itemize}

    \medskip
    \textbf{This task: all attention layers \cmark{}}\\
    Compare to 1-hop chain where L3--L9 have\\
    thick teal U-arcs (attn silent for 7 layers).

    \medskip
    \textbf{L11 final attention}: $r=0.99$ — the highest\\
    observed across all tasks.  The last-layer\\
    integrator works \emph{harder} for deeper chains.

    \medskip
    \textbf{L2 FFN spike}: $r=1.12$ — still above stream\\
    magnitude, but lower than 1-hop ($r=1.84$).\\
    Early FFN pre-processing budget shrinks\\
    as the task distributes computation deeper.
  };

\end{tikzpicture}
\end{frame}


% ── Image slide 2: GPT-2 Counterfactual ──────────────────────────────────
% Roughly square aspect ratio — shows the logic-type routing pattern.
% Key feature: L4 is the only silent-attn layer; L11 attn r=0.93.
\begin{frame}[plain]
\begin{tikzpicture}[remember picture, overlay]

  \fill[attnred] (current page.north west) rectangle
                 ([yshift=-0.9cm]current page.north east);
  \node[anchor=west, text=white, font=\bfseries\footnotesize]
        at ([xshift=0.4cm, yshift=-0.45cm]current page.north west)
        {GPT-2 (12L$\times$12H) \quad Counterfactual \quad $k=20$ \quad
         \texttt{compare\_gpt2\_task5.png}};

  \node[anchor=north west] (img)
        at ([xshift=0.2cm,yshift=-1.0cm]current page.north west) {%
    \includegraphics[height=0.86\paperheight,
                     keepaspectratio]{compare_gpt2_task5.png}%
  };

  \node[anchor=north west, text width=6.8cm,
        font=\footnotesize, align=left]
        at ([xshift=5.5cm,yshift=-1.1cm]current page.north west) {

    \textcolor{attnred}{\textbf{Counterfactual task}}\\
    \emph{``In a world where cats bark and dogs meow,\\
    if you hear barking outside you think it is a''}

    \medskip
    \textbf{Silence: only L4 ($r=0.05$, borderline)}\\
    Almost every attention layer fires — this task\\
    requires the most distributed reasoning of\\
    the 6 tasks tested in GPT-2.

    \medskip
    \textbf{L10 FFN: $r=0.30$} — notably high compared\\
    to chain tasks ($r\approx 0.20$--$0.24$).\\
    The FFN gradient from L3 upward is steeper,\\
    suggesting more sustained semantic refinement.

    \medskip
    \textbf{Clusters with: syllogism} (both wake at L4/L5).\\
    Contrast: 1-hop chain wakes attention at L10;\\
    counterfactual has attention alive throughout.

    \medskip
    \textbf{k=20} (maximum for GPT-2, same as syllogism\\
    and negation) — attribution spread across the\\
    full active subgraph, not concentrated.
  };

\end{tikzpicture}
\end{frame}


% ── Image slide 3: GPT-2-medium 3-hop chain ──────────────────────────────
% Portrait aspect ratio — the only GPT-2-medium graph with no silence zone.
% Key feature: bilateral dead zone in 1-hop is ABSENT here.
\begin{frame}[plain]
\begin{tikzpicture}[remember picture, overlay]

  \fill[mlpblue] (current page.north west) rectangle
                 ([yshift=-0.9cm]current page.north east);
  \node[anchor=west, text=white, font=\bfseries\footnotesize]
        at ([xshift=0.4cm, yshift=-0.45cm]current page.north west)
        {GPT-2-medium (24L$\times$16H) \quad 3-hop Chain \quad $k=36$ \quad
         \texttt{compare\_gpt2\_medium\_task2.png}};

  \node[anchor=north west] (img)
        at ([xshift=0.2cm,yshift=-1.0cm]current page.north west) {%
    \includegraphics[height=0.86\paperheight,
                     keepaspectratio]{compare_gpt2_medium_task2.png}%
  };

  \node[anchor=north west, text width=7.5cm,
        font=\footnotesize, align=left]
        at ([xshift=4.2cm,yshift=-1.1cm]current page.north west) {

    \textcolor{mlpblue}{\textbf{The anomalous task in GPT-2-medium}}

    \medskip
    \textbf{Every layer shows attn \cmark{}}\\
    All other 5 tasks have long silence zones\\
    (up to 17 consecutive attn-inactive layers).\\
    The 3-hop chain is categorically different.

    \medskip
    \textbf{Compare with 1-hop (not shown):}
    \begin{itemize}\setlength\itemsep{1pt}
      \item 1-hop: attn silent L4--L20; \emph{both} attn\\
            and FFN silent at L7--L13
      \item 3-hop: \textbf{zero silence} anywhere
    \end{itemize}

    \medskip
    \textbf{L0 FFN: $r=36.3$} (same huge spike as other\\
    tasks — the L0 FFN rewrites the token regardless\\
    of task type).

    \medskip
    \textbf{L23 attn: $r=0.42$} — universal output integrator\\
    active across all GPT-2-medium tasks.\\[3pt]
    \textbf{L22 FFN: $r=0.29$} — pre-final FFN spike,\\
    consistent across all tasks.

    \medskip
    \textbf{Open question}: Is the activation of all 24\\
    attention layers a sign of capacity saturation\\
    (model cannot solve 3-hop without every layer)?\\
    Or is it a qualitatively different circuit?
  };

\end{tikzpicture}
\end{frame}

% ─────────────────────────────────────────────────────────────────────────
\end{document}
